\documentclass[12pt, ukrainian]{article}
\usepackage[a4paper]{geometry}
\setlength\parindent{0mm}
\geometry{top=1.1cm,bottom=1.1cm,left=1.0cm,right=1.0cm}
\pagestyle{empty}
\usepackage[T2A]{fontenc}
\usepackage[utf8]{inputenc}
\usepackage[ukrainian]{babel}
\usepackage{graphicx}
\usepackage{caption}
\DeclareCaptionFormat{nocaption}{}
\captionsetup{format=nocaption,aboveskip=0pt,belowskip=0pt}
\usepackage[Export]{adjustbox} 
\adjustboxset{max size={0.9\linewidth}{0.9\paperheight}}
\begin{document}
%begin
\begin {Huge}
{{04.12.2023 Контрольна робота \No 2, варіант  \No \textbf { 101 }\par}}

\begin{center}Частина 2
\end{center}
\vspace{10mm}
\end {Huge}

{
\begin {Large}
Якість коду та економність обчислень оцінюється по обом завданням.

Співбесіда за роботою можлива.

\vspace{5mm}

%beginitem
1. Написати функцію, що знаходить найдовший відрізок списку цілих, у десятковому запису чисел якого є цифра 7. Повертає довжину знайденого відрізку або $0$ (за відсутності). Використовувати додаткові структури даних (у тому числі рядки) заборонено.

\vspace{10mm}

%enditem


%beginitem
2. 
Написати функцію, що повертає \textbf{новий список} з записаною в нього послідовністю чисел
$$\scalebox{1.3}{$(-2)^{k}\frac{(2n+k)!}{(2k+1)!(2n-k)!}$}, \quad k=1,\ldots,n-2$$
Оцінюється економність обчислень. Використовувати **, вкладені цикли та виклики функцій \textbf{заборонено}.

%enditem
\end{Large}\newpage

%end



%begin
\begin {Huge}
{{04.12.2023 Контрольна робота \No 2, варіант  \No \textbf { 102 }\par}}

\begin{center}Частина 2
\end{center}
\vspace{10mm}
\end {Huge}

{
\begin {Large}
Якість коду та економність обчислень оцінюється по обом завданням.

Співбесіда за роботою можлива.

\vspace{5mm}

%beginitem
1. Написати функцію, що для списку цілих знаходить найдовший відрізок, в якому всі сусідні числа відрізняються якнайменш на 3. Якщо таких відрізків кілька, то повернути перший (як індекс початку та довжину). Використовувати додаткові структури даних (у тому числі рядки) заборонено.

\vspace{10mm}

%enditem


%beginitem
2. 
Написати функцію, що повертає \textbf{новий список} з записаною в нього послідовністю чисел
$$\scalebox{1.2}{$C_{n+2k}^{k+1}$}, \quad k=0,1,\ldots,n$$
Оцінюється економність обчислень. Використовувати **, вкладені цикли та виклики функцій \textbf{заборонено}.

%enditem
\end{Large}\newpage

%end



%begin
\begin {Huge}
{{04.12.2023 Контрольна робота \No 2, варіант  \No \textbf { 103 }\par}}

\begin{center}Частина 2
\end{center}
\vspace{10mm}
\end {Huge}

{
\begin {Large}
Якість коду та економність обчислень оцінюється по обом завданням.

Співбесіда за роботою можлива.

\vspace{5mm}

%beginitem
1. Написати функцію, що знаходить найбільше число з списку цілих, у десятковому запису якого нема цифри 4. Повертає індекс першого входження цього числа або $-1$ (за відсутності). Використовувати додаткові структури даних (у тому числі рядки) заборонено.

\vspace{10mm}

%enditem


%beginitem
2. 
Написати функцію, що повертає \textbf{новий список} з записаною в нього послідовністю чисел
$$\scalebox{1.2}{$C_{2n+k+1}^{n+k}$}, \quad k=1,\ldots,n$$
Оцінюється економність обчислень. Використовувати **, вкладені цикли та виклики функцій \textbf{заборонено}.

%enditem
\end{Large}\newpage

%end



%begin
\begin {Huge}
{{04.12.2023 Контрольна робота \No 2, варіант  \No \textbf { 104 }\par}}

\begin{center}Частина 2
\end{center}
\vspace{10mm}
\end {Huge}

{
\begin {Large}
Якість коду та економність обчислень оцінюється по обом завданням.

Співбесіда за роботою можлива.

\vspace{5mm}

%beginitem
1. Написати функцію, що для списку цілих знаходить довжину найдовшого відрізку, в якому всі сусідні числа відрізняються якнайменш на 2. Використовувати додаткові структури даних (у тому числі рядки) заборонено.

\vspace{10mm}

%enditem


%beginitem
2. 
Написати функцію, що повертає \textbf{новий список} з записаною в нього послідовністю чисел
$$\scalebox{1.2}{$C_{3k-1}^{k}$}, \quad k=2,\ldots,n$$
Оцінюється економність обчислень. Використовувати **, вкладені цикли та виклики функцій \textbf{заборонено}.

%enditem
\end{Large}\newpage

%end



%begin
\begin {Huge}
{{04.12.2023 Контрольна робота \No 2, варіант  \No \textbf { 105 }\par}}

\begin{center}Частина 2
\end{center}
\vspace{10mm}
\end {Huge}

{
\begin {Large}
Якість коду та економність обчислень оцінюється по обом завданням.

Співбесіда за роботою можлива.

\vspace{5mm}

%beginitem
1. Написати функцію, що для списку цілих знаходить найдовший відрізок, в якому всі сусідні числа різні. Якщо таких відрізків кілька, то повернути останній (як індекс початку та довжину). Використовувати додаткові структури даних (у тому числі рядки) заборонено.

\vspace{10mm}

%enditem


%beginitem
2. 
Написати функцію, що повертає \textbf{новий список} з записаною в нього послідовністю чисел
$$\scalebox{1.2}{$(-1)^{k}C_{3k-1}^{k}$}, \quad k=2,\ldots,n$$
Оцінюється економність обчислень. Використовувати **, вкладені цикли та виклики функцій \textbf{заборонено}.

%enditem
\end{Large}\newpage

%end



%begin
\begin {Huge}
{{04.12.2023 Контрольна робота \No 2, варіант  \No \textbf { 106 }\par}}

\begin{center}Частина 2
\end{center}
\vspace{10mm}
\end {Huge}

{
\begin {Large}
Якість коду та економність обчислень оцінюється по обом завданням.

Співбесіда за роботою можлива.

\vspace{5mm}

%beginitem
1. Написати функцію, що повертає довжину найдовшого відрізку послідовності, що складається з чисел, у сімковому запису якого нема цифри 2. За відсутності має повертати $0$. Використовувати додаткові структури даних (у тому числі рядки) заборонено.

\vspace{10mm}

%enditem


%beginitem
2. 
Написати функцію, що повертає \textbf{новий список} з записаною в нього послідовністю чисел
$$\scalebox{1.2}{$(-1)^{k}C_{2n+k+1}^{n+k}$}, \quad k=0,1,\ldots,n$$
Оцінюється економність обчислень. Використовувати **, вкладені цикли та виклики функцій \textbf{заборонено}.

%enditem
\end{Large}\newpage

%end



%begin
\begin {Huge}
{{04.12.2023 Контрольна робота \No 2, варіант  \No \textbf { 107 }\par}}

\begin{center}Частина 2
\end{center}
\vspace{10mm}
\end {Huge}

{
\begin {Large}
Якість коду та економність обчислень оцінюється по обом завданням.

Співбесіда за роботою можлива.

\vspace{5mm}

%beginitem
1. Написати функцію, що знаходить найдовший відрізок списку цілих, у десятковому запису чисел якого є цифра 7. Повертає довжину знайденого відрізку або $0$ (за відсутності). Використовувати додаткові структури даних (у тому числі рядки) заборонено.

\vspace{10mm}

%enditem


%beginitem
2. 
Написати функцію, що повертає \textbf{новий список} з записаною в нього послідовністю чисел
$$\scalebox{1.3}{$(-3)^{k}\frac{(3k)!!(3n+k)!}{(3k+1)!(3n-k)!}$}, \quad k=1,\ldots,n$$
Оцінюється економність обчислень. Використовувати **, вкладені цикли та виклики функцій \textbf{заборонено}.

%enditem
\end{Large}\newpage

%end



%begin
\begin {Huge}
{{04.12.2023 Контрольна робота \No 2, варіант  \No \textbf { 108 }\par}}

\begin{center}Частина 2
\end{center}
\vspace{10mm}
\end {Huge}

{
\begin {Large}
Якість коду та економність обчислень оцінюється по обом завданням.

Співбесіда за роботою можлива.

\vspace{5mm}

%beginitem
1. Написати функцію, що повертає число з списку цілих, що має найменшу кількість різних простих дільників. Якщо таких чисел кілька, то повернути найбільше. Використовувати додаткові структури даних (у тому числі рядки) заборонено.

\vspace{10mm}

%enditem


%beginitem
2. 
Написати функцію, що повертає \textbf{новий список} з записаною в нього послідовністю чисел
$$\scalebox{1.2}{$\frac{(2k+1)!!}{(k!)^2}$}, \quad k=1,\ldots,n$$
Оцінюється економність обчислень. Використовувати **, вкладені цикли та виклики функцій \textbf{заборонено}.

%enditem
\end{Large}\newpage

%end



%begin
\begin {Huge}
{{04.12.2023 Контрольна робота \No 2, варіант  \No \textbf { 109 }\par}}

\begin{center}Частина 2
\end{center}
\vspace{10mm}
\end {Huge}

{
\begin {Large}
Якість коду та економність обчислень оцінюється по обом завданням.

Співбесіда за роботою можлива.

\vspace{5mm}

%beginitem
1. Написати функцію, що знаходить найбільше число з списку цілих, у десятковому запису якого нема цифри 3. Повертає індекс останнього входження цього числа або $-1$ (за відсутності). Використовувати додаткові структури даних (у тому числі рядки) заборонено.

\vspace{10mm}

%enditem


%beginitem
2. 
Написати функцію, що повертає \textbf{новий список} з записаною в нього послідовністю чисел
$$\scalebox{1.2}{$\frac{(2k+1)!!}{(k+1)!(k-1)!}$}, \quad k=2,\ldots,n$$
Оцінюється економність обчислень. Використовувати **, вкладені цикли та виклики функцій \textbf{заборонено}.

%enditem
\end{Large}\newpage

%end



%begin
\begin {Huge}
{{04.12.2023 Контрольна робота \No 2, варіант  \No \textbf { 110 }\par}}

\begin{center}Частина 2
\end{center}
\vspace{10mm}
\end {Huge}

{
\begin {Large}
Якість коду та економність обчислень оцінюється по обом завданням.

Співбесіда за роботою можлива.

\vspace{5mm}

%beginitem
1. Написати функцію, що повертає найбільше число з списку цілих, що має не більше трьох різних простих дільників. За відсутності має повертати None. Використовувати додаткові структури даних (у тому числі рядки) заборонено.

\vspace{10mm}

%enditem


%beginitem
2. 
Написати функцію, що повертає \textbf{новий список} з записаною в нього послідовністю чисел
$$\scalebox{1.2}{$(-1)^{k}C_{n+k}^{k}$}, \quad k=0,1,\ldots,n$$
Оцінюється економність обчислень. Використовувати **, вкладені цикли та виклики функцій \textbf{заборонено}.

%enditem
\end{Large}\newpage

%end



%begin
\begin {Huge}
{{04.12.2023 Контрольна робота \No 2, варіант  \No \textbf { 111 }\par}}

\begin{center}Частина 2
\end{center}
\vspace{10mm}
\end {Huge}

{
\begin {Large}
Якість коду та економність обчислень оцінюється по обом завданням.

Співбесіда за роботою можлива.

\vspace{5mm}

%beginitem
1. Написати функцію, що повертає найменше число з списку цілих, що має не більше трьох різних простих дільників. За відсутності має повертати None. Використовувати додаткові структури даних (у тому числі рядки) заборонено.

\vspace{10mm}

%enditem


%beginitem
2. 
Написати функцію, що повертає \textbf{новий список} з записаною в нього послідовністю чисел
$$\scalebox{1.2}{$C_{2n+k+1}^{n+k}$}, \quad k=0,1,\ldots,n$$
Оцінюється економність обчислень. Використовувати **, вкладені цикли та виклики функцій \textbf{заборонено}.

%enditem
\end{Large}\newpage

%end



%begin
\begin {Huge}
{{04.12.2023 Контрольна робота \No 2, варіант  \No \textbf { 112 }\par}}

\begin{center}Частина 2
\end{center}
\vspace{10mm}
\end {Huge}

{
\begin {Large}
Якість коду та економність обчислень оцінюється по обом завданням.

Співбесіда за роботою можлива.

\vspace{5mm}

%beginitem
1. Написати функцію, що повертає число з списку цілих, що має парну кількість різних простих дільників. Якщо таких чисел кілька, то повернути найбільше. Використовувати додаткові структури даних (у тому числі рядки) заборонено.

\vspace{10mm}

%enditem


%beginitem
2. 
Написати функцію, що повертає \textbf{новий список} з записаною в нього послідовністю чисел
$$\scalebox{1.2}{$C_{n+k}^{2k}$}, \quad k=1,\ldots,n$$
Оцінюється економність обчислень. Використовувати **, вкладені цикли та виклики функцій \textbf{заборонено}.

%enditem
\end{Large}\newpage

%end



%begin
\begin {Huge}
{{04.12.2023 Контрольна робота \No 2, варіант  \No \textbf { 113 }\par}}

\begin{center}Частина 2
\end{center}
\vspace{10mm}
\end {Huge}

{
\begin {Large}
Якість коду та економність обчислень оцінюється по обом завданням.

Співбесіда за роботою можлива.

\vspace{5mm}

%beginitem
1. Написати функцію, що повертає число з списку цілих, що має найбільшу кількість різних простих дільників. Якщо таких чисел кілька, то повернути найменше. Використовувати додаткові структури даних (у тому числі рядки) заборонено.

\vspace{10mm}

%enditem


%beginitem
2. 
Написати функцію, що повертає \textbf{новий список} з записаною в нього послідовністю чисел
$$\scalebox{1.2}{$(-1)^{k}\frac{(2k+1)!!}{(k!)^2}$}, \quad k=1,\ldots,n$$
Оцінюється економність обчислень. Використовувати **, вкладені цикли та виклики функцій \textbf{заборонено}.

%enditem
\end{Large}\newpage

%end



%begin
\begin {Huge}
{{04.12.2023 Контрольна робота \No 2, варіант  \No \textbf { 114 }\par}}

\begin{center}Частина 2
\end{center}
\vspace{10mm}
\end {Huge}

{
\begin {Large}
Якість коду та економність обчислень оцінюється по обом завданням.

Співбесіда за роботою можлива.

\vspace{5mm}

%beginitem
1. Написати функцію, що знаходить найбільше число з списку цілих, у десятковому запису якого нема цифри 4. Повертає індекс першого входження цього числа або $-1$ (за відсутності). Використовувати додаткові структури даних (у тому числі рядки) заборонено.

\vspace{10mm}

%enditem


%beginitem
2. 
Написати функцію, що повертає \textbf{новий список} з записаною в нього послідовністю чисел
$$\scalebox{1.2}{$(-1)^{k}\frac{(2k+1)!!}{(k+1)!(k-1)!}$}, \quad k=2,\ldots,n$$
Оцінюється економність обчислень. Використовувати **, вкладені цикли та виклики функцій \textbf{заборонено}.

%enditem
\end{Large}\newpage

%end



%begin
\begin {Huge}
{{04.12.2023 Контрольна робота \No 2, варіант  \No \textbf { 115 }\par}}

\begin{center}Частина 2
\end{center}
\vspace{10mm}
\end {Huge}

{
\begin {Large}
Якість коду та економність обчислень оцінюється по обом завданням.

Співбесіда за роботою можлива.

\vspace{5mm}

%beginitem
1. Написати функцію, що повертає число з списку цілих, що має парну кількість різних простих дільників. Якщо таких чисел кілька, то повернути найбільше. Використовувати додаткові структури даних (у тому числі рядки) заборонено.

\vspace{10mm}

%enditem


%beginitem
2. 
Написати функцію, що повертає \textbf{новий список} з записаною в нього послідовністю чисел
$$\scalebox{1.2}{$C_{n+k}^{2k}$}, \quad k=0,1,\ldots,n$$
Оцінюється економність обчислень. Використовувати **, вкладені цикли та виклики функцій \textbf{заборонено}.

%enditem
\end{Large}\newpage

%end



%begin
\begin {Huge}
{{04.12.2023 Контрольна робота \No 2, варіант  \No \textbf { 116 }\par}}

\begin{center}Частина 2
\end{center}
\vspace{10mm}
\end {Huge}

{
\begin {Large}
Якість коду та економність обчислень оцінюється по обом завданням.

Співбесіда за роботою можлива.

\vspace{5mm}

%beginitem
1. Написати функцію, що знаходить найбільше число з списку цілих, у десятковому запису якого нема цифри 3. Повертає індекс останнього входження цього числа або $-1$ (за відсутності). Використовувати додаткові структури даних (у тому числі рядки) заборонено.

\vspace{10mm}

%enditem


%beginitem
2. 
Написати функцію, що повертає \textbf{новий список} з записаною в нього послідовністю чисел
$$\scalebox{1.2}{$C_{3k}^{k}$}, \quad k=2,\ldots,n$$
Оцінюється економність обчислень. Використовувати **, вкладені цикли та виклики функцій \textbf{заборонено}.

%enditem
\end{Large}\newpage

%end



%begin
\begin {Huge}
{{04.12.2023 Контрольна робота \No 2, варіант  \No \textbf { 117 }\par}}

\begin{center}Частина 2
\end{center}
\vspace{10mm}
\end {Huge}

{
\begin {Large}
Якість коду та економність обчислень оцінюється по обом завданням.

Співбесіда за роботою можлива.

\vspace{5mm}

%beginitem
1. Написати функцію, що для списку цілих знаходить найдовший відрізок, в якому всі сусідні числа відрізняються рівно на 2. Якщо таких відрізків кілька, то повернути останній (як індекс початку та довжину). Використовувати додаткові структури даних (у тому числі рядки) заборонено.

\vspace{10mm}

%enditem


%beginitem
2. 
Написати функцію, що повертає \textbf{новий список} з записаною в нього послідовністю чисел
$$\scalebox{1.3}{$(-2)^{k}\frac{(2n+k)!}{(2k+3)!(2n-k+1)!}$}, \quad k=1,\ldots,n-3$$
Оцінюється економність обчислень. Використовувати **, вкладені цикли та виклики функцій \textbf{заборонено}.

%enditem
\end{Large}\newpage

%end



%begin
\begin {Huge}
{{04.12.2023 Контрольна робота \No 2, варіант  \No \textbf { 118 }\par}}

\begin{center}Частина 2
\end{center}
\vspace{10mm}
\end {Huge}

{
\begin {Large}
Якість коду та економність обчислень оцінюється по обом завданням.

Співбесіда за роботою можлива.

\vspace{5mm}

%beginitem
1. Написати функцію, що знаходить найменше число з списку цілих, у десятковому запису якого нема цифри 5. Повертає індекс останнього входження цього числа або $-1$ (за відсутності). Використовувати додаткові структури даних (у тому числі рядки) заборонено.

\vspace{10mm}

%enditem


%beginitem
2. 
Написати функцію, що повертає \textbf{новий список} з записаною в нього послідовністю чисел
$$\scalebox{1.2}{$C_{2k-1}^{k}$}, \quad k=2,\ldots,n$$
Оцінюється економність обчислень. Використовувати **, вкладені цикли та виклики функцій \textbf{заборонено}.

%enditem
\end{Large}\newpage

%end



%begin
\begin {Huge}
{{04.12.2023 Контрольна робота \No 2, варіант  \No \textbf { 119 }\par}}

\begin{center}Частина 2
\end{center}
\vspace{10mm}
\end {Huge}

{
\begin {Large}
Якість коду та економність обчислень оцінюється по обом завданням.

Співбесіда за роботою можлива.

\vspace{5mm}

%beginitem
1. Написати функцію, що повертає довжину найдовшого відрізку послідовності, що складається з чисел, у сімковому запису якого нема цифри 2. За відсутності має повертати $0$. Використовувати додаткові структури даних (у тому числі рядки) заборонено.

\vspace{10mm}

%enditem


%beginitem
2. 
Написати функцію, що повертає \textbf{новий список} з записаною в нього послідовністю чисел
$$\scalebox{1.2}{$(-1)^{k}C_{n+2k}^{k}$}, \quad k=0,1,\ldots,n-2$$
Оцінюється економність обчислень. Використовувати **, вкладені цикли та виклики функцій \textbf{заборонено}.

%enditem
\end{Large}\newpage

%end



%begin
\begin {Huge}
{{04.12.2023 Контрольна робота \No 2, варіант  \No \textbf { 120 }\par}}

\begin{center}Частина 2
\end{center}
\vspace{10mm}
\end {Huge}

{
\begin {Large}
Якість коду та економність обчислень оцінюється по обом завданням.

Співбесіда за роботою можлива.

\vspace{5mm}

%beginitem
1. Написати функцію, що повертає найменше число з списку цілих, що має не більше трьох різних простих дільників. За відсутності має повертати None. Використовувати додаткові структури даних (у тому числі рядки) заборонено.

\vspace{10mm}

%enditem


%beginitem
2. 
Написати функцію, що повертає \textbf{новий список} з записаною в нього послідовністю чисел
$$\scalebox{1.2}{$(-1)^{k}\frac{(2k)!!}{(k!)^2}$}, \quad k=1,\ldots,n$$
Оцінюється економність обчислень. Використовувати **, вкладені цикли та виклики функцій \textbf{заборонено}.

%enditem
\end{Large}\newpage

%end



%begin
\begin {Huge}
{{04.12.2023 Контрольна робота \No 2, варіант  \No \textbf { 121 }\par}}

\begin{center}Частина 2
\end{center}
\vspace{10mm}
\end {Huge}

{
\begin {Large}
Якість коду та економність обчислень оцінюється по обом завданням.

Співбесіда за роботою можлива.

\vspace{5mm}

%beginitem
1. Написати функцію, що повертає число з списку цілих, що має найбільшу кількість різних простих дільників. Якщо таких чисел кілька, то повернути найменше. Використовувати додаткові структури даних (у тому числі рядки) заборонено.

\vspace{10mm}

%enditem


%beginitem
2. 
Написати функцію, що повертає \textbf{новий список} з записаною в нього послідовністю чисел
$$\scalebox{1.2}{$C_{n+k}^{k}$}, \quad k=0,1,\ldots,n-2$$
Оцінюється економність обчислень. Використовувати **, вкладені цикли та виклики функцій \textbf{заборонено}.

%enditem
\end{Large}\newpage

%end



%begin
\begin {Huge}
{{04.12.2023 Контрольна робота \No 2, варіант  \No \textbf { 122 }\par}}

\begin{center}Частина 2
\end{center}
\vspace{10mm}
\end {Huge}

{
\begin {Large}
Якість коду та економність обчислень оцінюється по обом завданням.

Співбесіда за роботою можлива.

\vspace{5mm}

%beginitem
1. Написати функцію, що повертає найбільше число з списку цілих, що має не більше трьох різних простих дільників. За відсутності має повертати None. Використовувати додаткові структури даних (у тому числі рядки) заборонено.

\vspace{10mm}

%enditem


%beginitem
2. 
Написати функцію, що повертає \textbf{новий список} з записаною в нього послідовністю чисел
$$\scalebox{1.2}{$C_{n+2k}^{k}$}, \quad k=1,\ldots,n$$
Оцінюється економність обчислень. Використовувати **, вкладені цикли та виклики функцій \textbf{заборонено}.

%enditem
\end{Large}\newpage

%end



%begin
\begin {Huge}
{{04.12.2023 Контрольна робота \No 2, варіант  \No \textbf { 123 }\par}}

\begin{center}Частина 2
\end{center}
\vspace{10mm}
\end {Huge}

{
\begin {Large}
Якість коду та економність обчислень оцінюється по обом завданням.

Співбесіда за роботою можлива.

\vspace{5mm}

%beginitem
1. Написати функцію, що для списку цілих знаходить найдовший відрізок, в якому всі сусідні числа відрізняються якнайменш на 3. Якщо таких відрізків кілька, то повернути перший (як індекс початку та довжину). Використовувати додаткові структури даних (у тому числі рядки) заборонено.

\vspace{10mm}

%enditem


%beginitem
2. 
Написати функцію, що повертає \textbf{новий список} з записаною в нього послідовністю чисел
$$\scalebox{1.2}{$(-1)^{k}C_{n+k}^{2k}$}, \quad k=1,\ldots,n$$
Оцінюється економність обчислень. Використовувати **, вкладені цикли та виклики функцій \textbf{заборонено}.

%enditem
\end{Large}\newpage

%end



%begin
\begin {Huge}
{{04.12.2023 Контрольна робота \No 2, варіант  \No \textbf { 124 }\par}}

\begin{center}Частина 2
\end{center}
\vspace{10mm}
\end {Huge}

{
\begin {Large}
Якість коду та економність обчислень оцінюється по обом завданням.

Співбесіда за роботою можлива.

\vspace{5mm}

%beginitem
1. Написати функцію, що знаходить друге за абсолютною величиною число з списку цілих, у десятковому запису якого є цифра 5 або цифра 3. Повертає індекс останнього входження цього числа або $-1$ (за відсутності). Використовувати додаткові структури даних (у тому числі рядки) заборонено.

\vspace{10mm}

%enditem


%beginitem
2. 
Написати функцію, що повертає \textbf{новий список} з записаною в нього послідовністю чисел
$$\scalebox{1.3}{$(-2)^{k}\frac{(3k+1)!!(2n+k)!}{(k+2)!(2n-k)!}$}, \quad k=1,\ldots,n+2$$
Оцінюється економність обчислень. Використовувати **, вкладені цикли та виклики функцій \textbf{заборонено}.

%enditem
\end{Large}\newpage

%end



%begin
\begin {Huge}
{{04.12.2023 Контрольна робота \No 2, варіант  \No \textbf { 125 }\par}}

\begin{center}Частина 2
\end{center}
\vspace{10mm}
\end {Huge}

{
\begin {Large}
Якість коду та економність обчислень оцінюється по обом завданням.

Співбесіда за роботою можлива.

\vspace{5mm}

%beginitem
1. Написати функцію, що повертає число з списку цілих, що має найменшу кількість різних простих дільників. Якщо таких чисел кілька, то повернути найбільше. Використовувати додаткові структури даних (у тому числі рядки) заборонено.

\vspace{10mm}

%enditem


%beginitem
2. 
Написати функцію, що повертає \textbf{новий список} з записаною в нього послідовністю чисел
$$\scalebox{1.2}{$C_{n+2k}^{k+1}$}, \quad k=1,\ldots,n$$
Оцінюється економність обчислень. Використовувати **, вкладені цикли та виклики функцій \textbf{заборонено}.

%enditem
\end{Large}\newpage

%end



%begin
\begin {Huge}
{{04.12.2023 Контрольна робота \No 2, варіант  \No \textbf { 126 }\par}}

\begin{center}Частина 2
\end{center}
\vspace{10mm}
\end {Huge}

{
\begin {Large}
Якість коду та економність обчислень оцінюється по обом завданням.

Співбесіда за роботою можлива.

\vspace{5mm}

%beginitem
1. Написати функцію, що для списку цілих знаходить довжину найдовшого відрізку, в якому всі сусідні числа відрізняються якнайменш на 2. Використовувати додаткові структури даних (у тому числі рядки) заборонено.

\vspace{10mm}

%enditem


%beginitem
2. 
Написати функцію, що повертає \textbf{новий список} з записаною в нього послідовністю чисел
$$\scalebox{1.2}{$(-1)^{k}C_{n+k}^{k+1}$}, \quad k=1,\ldots,n$$
Оцінюється економність обчислень. Використовувати **, вкладені цикли та виклики функцій \textbf{заборонено}.

%enditem
\end{Large}\newpage

%end



%begin
\begin {Huge}
{{04.12.2023 Контрольна робота \No 2, варіант  \No \textbf { 127 }\par}}

\begin{center}Частина 2
\end{center}
\vspace{10mm}
\end {Huge}

{
\begin {Large}
Якість коду та економність обчислень оцінюється по обом завданням.

Співбесіда за роботою можлива.

\vspace{5mm}

%beginitem
1. Написати функцію, що повертає число з списку цілих, що має найбільшу кількість різних простих дільників. Якщо таких чисел кілька, то повернути найменше. Використовувати додаткові структури даних (у тому числі рядки) заборонено.

\vspace{10mm}

%enditem


%beginitem
2. 
Написати функцію, що повертає \textbf{новий список} з записаною в нього послідовністю чисел
$$\scalebox{1.2}{$(-2)^{k}\frac{(n+k)!}{(2k-3)!(2n-k+1)!}$}, \quad k=3,\ldots,n+2$$
Оцінюється економність обчислень. Використовувати **, вкладені цикли та виклики функцій \textbf{заборонено}.

%enditem
\end{Large}\newpage

%end



%begin
\begin {Huge}
{{04.12.2023 Контрольна робота \No 2, варіант  \No \textbf { 128 }\par}}

\begin{center}Частина 2
\end{center}
\vspace{10mm}
\end {Huge}

{
\begin {Large}
Якість коду та економність обчислень оцінюється по обом завданням.

Співбесіда за роботою можлива.

\vspace{5mm}

%beginitem
1. Написати функцію, що повертає довжину найдовшого відрізку послідовності, що складається з чисел, у сімковому запису якого нема цифри 2. За відсутності має повертати $0$. Використовувати додаткові структури даних (у тому числі рядки) заборонено.

\vspace{10mm}

%enditem


%beginitem
2. 
Написати функцію, що повертає \textbf{новий список} з записаною в нього послідовністю чисел
$$\scalebox{1.2}{$(-1)^{k}C_{2n+k}^{n+k}$}, \quad k=0,1,\ldots,n$$
Оцінюється економність обчислень. Використовувати **, вкладені цикли та виклики функцій \textbf{заборонено}.

%enditem
\end{Large}\newpage

%end



%begin
\begin {Huge}
{{04.12.2023 Контрольна робота \No 2, варіант  \No \textbf { 129 }\par}}

\begin{center}Частина 2
\end{center}
\vspace{10mm}
\end {Huge}

{
\begin {Large}
Якість коду та економність обчислень оцінюється по обом завданням.

Співбесіда за роботою можлива.

\vspace{5mm}

%beginitem
1. Написати функцію, що для списку цілих знаходить найдовший відрізок, в якому всі сусідні числа відрізняються якнайменш на 3. Якщо таких відрізків кілька, то повернути перший (як індекс початку та довжину). Використовувати додаткові структури даних (у тому числі рядки) заборонено.

\vspace{10mm}

%enditem


%beginitem
2. 
Написати функцію, що повертає \textbf{новий список} з записаною в нього послідовністю чисел
$$\scalebox{1.2}{$C_{n+2k}^{k}$}, \quad k=0,1,\ldots,n-2$$
Оцінюється економність обчислень. Використовувати **, вкладені цикли та виклики функцій \textbf{заборонено}.

%enditem
\end{Large}\newpage

%end



%begin
\begin {Huge}
{{04.12.2023 Контрольна робота \No 2, варіант  \No \textbf { 130 }\par}}

\begin{center}Частина 2
\end{center}
\vspace{10mm}
\end {Huge}

{
\begin {Large}
Якість коду та економність обчислень оцінюється по обом завданням.

Співбесіда за роботою можлива.

\vspace{5mm}

%beginitem
1. Написати функцію, що повертає число з списку цілих, що має парну кількість різних простих дільників. Якщо таких чисел кілька, то повернути найбільше. Використовувати додаткові структури даних (у тому числі рядки) заборонено.

\vspace{10mm}

%enditem


%beginitem
2. 
Написати функцію, що повертає \textbf{новий список} з записаною в нього послідовністю чисел
$$\scalebox{1.2}{$(-1)^{k}C_{n+k}^{k+1}$}, \quad k=0,1,\ldots,n$$
Оцінюється економність обчислень. Використовувати **, вкладені цикли та виклики функцій \textbf{заборонено}.

%enditem
\end{Large}\newpage

%end



%begin
\begin {Huge}
{{04.12.2023 Контрольна робота \No 2, варіант  \No \textbf { 131 }\par}}

\begin{center}Частина 2
\end{center}
\vspace{10mm}
\end {Huge}

{
\begin {Large}
Якість коду та економність обчислень оцінюється по обом завданням.

Співбесіда за роботою можлива.

\vspace{5mm}

%beginitem
1. Написати функцію, що для списку цілих знаходить довжину найдовшого відрізку, в якому всі сусідні числа відрізняються якнайменш на 2. Використовувати додаткові структури даних (у тому числі рядки) заборонено.

\vspace{10mm}

%enditem


%beginitem
2. 
Написати функцію, що повертає \textbf{новий список} з записаною в нього послідовністю чисел
$$\scalebox{1.2}{$(-1)^{k}C_{2k+1}^{k}$}, \quad k=2,\ldots,n$$
Оцінюється економність обчислень. Використовувати **, вкладені цикли та виклики функцій \textbf{заборонено}.

%enditem
\end{Large}\newpage

%end



%begin
\begin {Huge}
{{04.12.2023 Контрольна робота \No 2, варіант  \No \textbf { 132 }\par}}

\begin{center}Частина 2
\end{center}
\vspace{10mm}
\end {Huge}

{
\begin {Large}
Якість коду та економність обчислень оцінюється по обом завданням.

Співбесіда за роботою можлива.

\vspace{5mm}

%beginitem
1. Написати функцію, що повертає число з списку цілих, що має найменшу кількість різних простих дільників. Якщо таких чисел кілька, то повернути найбільше. Використовувати додаткові структури даних (у тому числі рядки) заборонено.

\vspace{10mm}

%enditem


%beginitem
2. 
Написати функцію, що повертає \textbf{новий список} з записаною в нього послідовністю чисел
$$\scalebox{1.2}{$C_{2k+1}^{k}$}, \quad k=2,\ldots,n$$
Оцінюється економність обчислень. Використовувати **, вкладені цикли та виклики функцій \textbf{заборонено}.

%enditem
\end{Large}\newpage

%end



%begin
\begin {Huge}
{{04.12.2023 Контрольна робота \No 2, варіант  \No \textbf { 133 }\par}}

\begin{center}Частина 2
\end{center}
\vspace{10mm}
\end {Huge}

{
\begin {Large}
Якість коду та економність обчислень оцінюється по обом завданням.

Співбесіда за роботою можлива.

\vspace{5mm}

%beginitem
1. Написати функцію, що знаходить найменше число з списку цілих, у десятковому запису якого нема цифри 5. Повертає індекс останнього входження цього числа або $-1$ (за відсутності). Використовувати додаткові структури даних (у тому числі рядки) заборонено.

\vspace{10mm}

%enditem


%beginitem
2. 
Написати функцію, що повертає \textbf{новий список} з записаною в нього послідовністю чисел
$$\scalebox{1.2}{$(-1)^{k}C_{2k-1}^{k}$}, \quad k=2,\ldots,n$$
Оцінюється економність обчислень. Використовувати **, вкладені цикли та виклики функцій \textbf{заборонено}.

%enditem
\end{Large}\newpage

%end



%begin
\begin {Huge}
{{04.12.2023 Контрольна робота \No 2, варіант  \No \textbf { 134 }\par}}

\begin{center}Частина 2
\end{center}
\vspace{10mm}
\end {Huge}

{
\begin {Large}
Якість коду та економність обчислень оцінюється по обом завданням.

Співбесіда за роботою можлива.

\vspace{5mm}

%beginitem
1. Написати функцію, що знаходить найбільше число з списку цілих, у десятковому запису якого нема цифри 3. Повертає індекс останнього входження цього числа або $-1$ (за відсутності). Використовувати додаткові структури даних (у тому числі рядки) заборонено.

\vspace{10mm}

%enditem


%beginitem
2. 
Написати функцію, що повертає \textbf{новий список} з записаною в нього послідовністю чисел
$$\scalebox{1.2}{$(-1)^{k}C_{n+2k}^{k+1}$}, \quad k=1,\ldots,n$$
Оцінюється економність обчислень. Використовувати **, вкладені цикли та виклики функцій \textbf{заборонено}.

%enditem
\end{Large}\newpage

%end



%begin
\begin {Huge}
{{04.12.2023 Контрольна робота \No 2, варіант  \No \textbf { 135 }\par}}

\begin{center}Частина 2
\end{center}
\vspace{10mm}
\end {Huge}

{
\begin {Large}
Якість коду та економність обчислень оцінюється по обом завданням.

Співбесіда за роботою можлива.

\vspace{5mm}

%beginitem
1. Написати функцію, що повертає найбільше число з списку цілих, що має не більше трьох різних простих дільників. За відсутності має повертати None. Використовувати додаткові структури даних (у тому числі рядки) заборонено.

\vspace{10mm}

%enditem


%beginitem
2. 
Написати функцію, що повертає \textbf{новий список} з записаною в нього послідовністю чисел
$$\scalebox{1.2}{$(-1)^{k}\frac{(2k)!!}{(k+1)!(k-1)!}$}, \quad k=2,\ldots,n$$
Оцінюється економність обчислень. Використовувати **, вкладені цикли та виклики функцій \textbf{заборонено}.

%enditem
\end{Large}\newpage

%end



\end{document}

